Les requêtes sont disponibles dans le notebook Jupyter suivant~: [TODO ajouter lien git backmatter]

\section*{Choix méthodologique}

\textit{Il a été choisi de réaliser ces requêtes via l'API de Wikidata pour~:}
\begin{enumerate}
	\item \textit{Stocker facilement les requêtes et pouvoir les modifier et réexécuter;}
	\item \textit{Pouvoir réaliser des boucles en Python et morceler certaines requêtes qui surchargeaient le SPARQL Endpoint;}
\end{enumerate}

\section*{Étape 1. Extraction des termes}

\subsection*{Solutions abandonnées}

\begin{enumerate}
	\item Requêter tous les termes liés à l'aéronautique en faisant une recherche plein texte (pages contenant \textit{aeronautics, aircraft, ...})~:
	\begin{itemize}
		\item La requête étant trop importante, le serveur \textit{timeout} systématiquement même en la limitant à un petit nombre de réponses (car il recherche l'intégralité des pages Wikidata pour trouver des correspondances).
	\end{itemize}
	
	\item Requêter tous les enfants des termes liés à l'aéronautique rencontrés (constructeurs, métiers, sciences, etc.)~:
	\begin{itemize}
		\item Prometteur, mais limité car~:
		\begin{itemize}
			\item On ne pense pas à toutes les possibilités ;
			\item La nature des relations que l'on pourrait considérer comme « enfant » sont trop différentes selon le type de terme récupéré (exemple~: les relations hiérarchiques internes à Wikidata très génériques comme « nature de l'élément » et « sous-classe de » sont faciles à récupérer pour construire un arbre, mais celui-ci est souvent trop général et peu pertinent pour un sujet lié à l'aéronautique. Pour certains termes comme les constructeurs, il faut rechercher par exemple quelle relation relie une filiale à une organisation, un modèle d'avion à son constructeur, etc. ; relations qui ne sont pas toujours évidentes ni uniformément utilisées selon les pages.) ;
			\item Certaines pages récupérées ne sont pas pertinentes.
		\end{itemize}
	\end{itemize}
\end{enumerate}

\subsection*{Solution retenue}

Pour obtenir un maximum de pages liées, le plus efficace a semblé être de faire une requête par les propriétés. La requête étant lourde, elle a été morcelée dans le \textit{notebook} Jupyter pour éviter de surcharger le serveur~:

\begin{enumerate}
	\item Recherche des propriétés dans lesquelles un terme donné est utilisé ;
	\item Récupération de toutes les pages liées à ce terme par une propriété ;
	\item Application de cette démarche à tous les termes génériques liés à l'aéronautique jugés pertinents.
\end{enumerate}

Cette méthode a ainsi permis de récupérer un maximum de termes directement reliés aux principaux concepts déjà présents dans les thésaurus aéronautiques du \mae\footnote{Les extractions réalisées sont disponibles dans le tableau suivant \refinterne{tab:wikidata}}.


\begin{table}[htbp]
	\centering
	\begin{tabular}{lll}
		\toprule
		\textbf{ID} & \textbf{Terme} & \textbf{Résultats}\\
		\midrule
		Q8421 & Aéronautique & 448 \\
		Q2876213 & Aerospace & 79  \\
		Q22719 & Astronautique & 358  \\
		Q765633 & Aviation & 3057  \\
		Q1757562 & Transport aérien & 581\\
		Q11436 & Aéronef & 10493 \\
		Q107009743 & Constructeur aéronautique & 43  \\
		Q936518 & Constructeur aérospatial & 926 \\
		Q3477363 & Industrie aérospatiale & 1734\\
		Q1875606 & Industrie aéronautique & 433 \\
		Q4055832 & Equipement aéronautique & 68 \\
		Q16693356 & Composant aéronautique & 629 \\
		Q17301966 & Aérostation & 6 \\
		Q206814 & Aviation civile & 427 \\
		Q627716 & Aviation militaire & 418 \\
		Q21152775 & Aviation occurrence & 2 \\
		Q744913 & Accident aérien & 4653 \\
		Q3002150 & Crash aérien & 1023  \\
		Q108284447 & Accident ou incident aérien & 186 \\
		Q3798668 & Aerospace engineering & 379  \\
		Q4828699 & Aeronautical engineering & 164 \\
		\midrule
		\textbf{Total} & \textbf{26107} & \\
		\bottomrule
	\end{tabular}
	\caption{Suivi des extractions Wikidata par propriétés}
	\label{tab:wikidata}
\end{table}

\section*{Étape 2. Nettoyage}

\begin{itemize}
	\item Suppression des lexèmes récupérés (catégorie particulière de pages Wikidata, non pertinente, ID commençant par L\ldots)
	\item Suppression des « Catégories » Wikidata
	\item Regroupement des termes parents
	\item Suppression des doublons
\end{itemize}

\section*{Étape 3. Enrichissement}

Récupération des parents (propriétés P31 et P279) des termes extraits, pour approfondir la hiérarchie.

\section*{Étape 4. Nettoyage final}

Suppression des doublons, suppression manuelle de termes non pertinents.

\section*{Tableau final}

\begin{itemize}
	\item 18 392 entrées comprenant les termes directement liés aux termes renseignés dans le tableau.
\end{itemize}

Structure~:
\begin{itemize}
	\item ID Wikidata
	\item Titre
	\item Description
	\item ID terme générique
	\item Terme générique
\end{itemize}

Autre tableau~: Liste des termes génériques.

\section*{Conclusion}

Cette approche, qui a semblé être la plus efficace pour récupérer un maximum de termes pertinents dans un temps limité, présente cependant ses limites~: 
\begin{itemize}
	\item son utilisation est complexe~: les tentatives pour récupérer une base de données pertinentes à partir d'un même modèle de requête se sont révélées trop lourdes, ou trop précises pour arriver à une base de donnée intéressante,
	\item cette extraction provenant de Wikidata -- qui est donc davantage une ontologie plutôt qu'un thésaurus -- produit une base de donnée très dispersée, avec une multitude de termes génériques qu'il aurait fallu rattacher artificiellement à des termes de tête pour arriver à un résultat facile à lire,
	\item elle montre enfin les limites de Wikidata -- qui pour ces sujets pointus présente des inégalités dans ses articles (choix différents de propriétés pour décrire une même relation, absence de titre d'article, ...).
\end{itemize}


